% ────────────────────────────────────────
%  ادامهٔ فصل ۳ — از جایی که قطع شد
% ────────────────────────────────────────

\begin{warningBox}[موضع گنجی — ادامه]
\begin{itemize}[label=\textcolor{cAccent}{$\bullet$}, rightmargin=2em]
    \item خمینی یک \keyword{فاشیست مذهبی} بود که قصد تأسیس دولت تمامیت‌خواه داشت.\footnote{اکبر گنجی، \textit{تلقی فاشیستی از دین و حکومت}، تهران: طرح نو، ۱۳۷۹، ص ۱۸–۳۵.}
    \item وعده‌های پاریس نه «ابهام ذهنی» بود و نه «تکامل فکری» — بلکه \keyword{دروغ محض} بود.
    \item فاصلهٔ میان «حکومت اسلامی» (۱۳۴۸) و «مقامی نخواهم داشت» (۱۳۵۷) فقط با فرض فریب آگاهانه قابل توضیح است.
    \item خمینی از ابتدا یک «پروژهٔ تمامیت‌خواه» داشت و تمام ابزارها — از دادگاه‌های انقلاب تا سپاه پاسداران تا حزب جمهوری اسلامی — بخشی از این پروژه بودند.
    \item سخن از «اجبار شرایط» خود-فریبی است: خمینی شرایط را \keyword{خلق} کرد، نه اینکه به آن‌ها \keyword{واکنش} نشان داد.\footnote{گنجی در مصاحبه‌های متأخرش (پس از ۱۳۸۸) از خمینی به‌عنوان «استالین ایرانی» یاد کرده است.}
\end{itemize}
\end{warningBox}

\subsection{هما کاتوزیان — نقد از زاویهٔ حقوقی-تاریخی}

کاتوزیان، اقتصاددان و مورخ ایرانی و استاد دانشگاه \en{Oxford}، بدون آنکه لزوماً واژهٔ «فریب» را به‌کار ببرد، نشان می‌دهد که فرایند تدوین قانون اساسی \keyword{مهندسی‌شده} بود:

\begin{noteBox}[تحلیل کاتوزیان دربارهٔ قانون اساسی]
\begin{enumerate}[label=\textcolor{cGold}{\bfseries\arabic*.}, rightmargin=2em]
    \item پیش‌نویس اولیهٔ قانون اساسی (تهیه‌شده توسط تیم حقوقدانان شامل دکتر حسن حبیبی و خود کاتوزیان و دیگران) یک سند \keyword{جمهوری دموکراتیک} بود — بدون ولایت فقیه.
    \item خمینی ابتدا این پیش‌نویس را تأیید کرد و حتی دستور داد به رفراندوم گذاشته شود.
    \item اما سپس — تحت فشار حزب جمهوری اسلامی به رهبری بهشتی — تصمیم به تشکیل «مجلس خبرگان بررسی قانون اساسی» گرفته شد.
    \item ترکیب مجلس خبرگان عمداً به‌گونه‌ای مهندسی شد که اکثریت از روحانیون تندرو باشند.
    \item نتیجه: حذف کامل ساختار جمهوری دموکراتیک و جایگزینی ولایت فقیه.\footnote{\en{Homa Katouzian, ``The Short-Lived Hopes of Iran's Secular Democrats,'' paper presented at various conferences; also see his contributions in Abrahamian et al.,} \textit{مقالات تاریخ معاصر ایران}، انتشارات مختلف.}
\end{enumerate}
\end{noteBox}

\subsection{ابوالحسن بنی‌صدر — شهادت قربانی}

بنی‌صدر، نخستین رئیس‌جمهور ایران (۱۳۵۹–۱۳۶۰) که پس از عزل به فرانسه گریخت، یکی از صریح‌ترین مدافعان تز دوم است. او استدلال می‌کند:

\begin{itemize}[label=\textcolor{cAccent}{$\blacktriangleright$}, rightmargin=2em]
    \item خمینی از \keyword{همان ماه‌های اول} پس از انقلاب، تمام تصمیمات کلیدی را شخصاً اتخاذ می‌کرد — از انتصابات تا سیاست خارجی.
    \item ادعای «بازگشت به قم» از ابتدا صوری بود: خمینی حتی در قم هم از طریق احمد خمینی و شورای انقلاب، مرکز تصمیم‌گیری باقی ماند.
    \item گروگان‌گیری سفارت (آبان ۱۳۵۸) — که به استعفای بازرگان و حذف لیبرال‌ها منجر شد — با \keyword{علم و آگاهی} حلقهٔ درونی خمینی (اگر نه خود او) سازمان‌دهی شد.
    \item عزل بنی‌صدر «پایان کار» ائتلاف بود: خمینی ثابت کرد که هر رئیس‌جمهوری‌ای که با ولایت فقیه تعارض داشته باشد قابل حذف است.\footnote{ابوالحسن بنی‌صدر، \textit{خیانت به امید}، پاریس: نشر انقلاب اسلامی، ۱۳۶۱، فصل‌های ۸–۱۲.}
\end{itemize}

\begin{warningBox}[تحفظ دربارهٔ شهادت بنی‌صدر]
باید در نظر داشت که بنی‌صدر خود \keyword{طرف دعواست}. او رقیب سیاسی حذف‌شدهٔ خمینی بود و انگیزهٔ قوی برای نشان‌دادن خمینی به‌عنوان «فریبکار» دارد. شهادت او ارزش مستند دارد اما نمی‌تواند به‌تنهایی اثبات‌کنندهٔ تز دوم باشد. مثلث‌سازی با منابع دیگر ضروری است.
\end{warningBox}

% ────────────────────────────────────────
\section{نقد تز دوم}
% ────────────────────────────────────────

\subsection{نقد از منظر تز اول (صداقت + اجبار)}

\begin{itemize}[label=\textcolor{cGreen}{$\checkmark$}, rightmargin=2em]
    \item \textbf{کتاب «حکومت اسلامی» اثبات فریب نیست:} ممکن است خمینی در ۱۳۴۸ نظریه‌ای داشته باشد اما در ۱۳۵۷ واقعاً تصور کند اجرای آن عملی نیست یا زمانش نرسیده. داشتن نظریه لزوماً به‌معنای قصد اجرای فوری نیست.
    \item \textbf{ثبات وعده‌ها:} اگر خمینی فریبکار بود، چرا این‌قدر \keyword{صریح} و \keyword{مکرر} وعده داد؟ فریبکار هوشمند ابهام می‌گذارد، نه وعدهٔ صریح.
    \item \textbf{انتصاب بازرگان:} اگر قصد تمرکز قدرت داشت، چرا یک لیبرال‌ غیرروحانی را نخست‌وزیر کرد؟ چرا از ابتدا یک روحانی تندرو نگذاشت؟
    \item \textbf{مسئلهٔ «توطئه‌اندیشی»:} تز دوم نیازمند فرض یک \keyword{طراحی بلندمدت} است. آیا واقعاً منطقی است که کسی ده سال نقشه بکشد و هیچ خطایی نکند؟ واقعیت معمولاً آشفته‌تر از توطئه‌هاست.\footnote{\en{This is the classical ``conspiracy theory'' critique. See Karl Popper, ``The Conspiracy Theory of Society'' in Conjectures and Refutations (1963).}}
\end{itemize}

\subsection{نقد از منظر تز سوم (تکامل)}

\begin{itemize}[label=\textcolor{cGold}{$\checkmark$}, rightmargin=2em]
    \item تز دوم فرض می‌کند خمینی از ابتدا \keyword{نقطهٔ نهایی} را می‌دانست. اما نظریهٔ ولایت فقیه خود تحولاتی داشت — از «کشف‌الاسرار» (۱۳۲۲) که فقط «نظارت» فقها را مطرح کرد، تا «حکومت اسلامی» (۱۳۴۸) که «ولایت» را مطرح کرد، تا «ولایت مطلقهٔ فقیه» (۱۳۶۶). این تحول مستمر با فرض «نقشهٔ از پیش آماده» سازگار نیست.
    \item «چندزبانگی» خمینی لزوماً فریب نبود — ممکن بود بازتاب \keyword{عدم قطعیت} فکری باشد. آدمی که خودش هنوز به نتیجهٔ قطعی نرسیده، در برابر مخاطبان مختلف حرف‌های مختلف می‌زند.
\end{itemize}

\subsection{نقد از منظر تز چهارم (فشار نخبگان)}

\begin{itemize}[label=\textcolor{cConsolid}{$\checkmark$}, rightmargin=2em]
    \item تز دوم همهٔ «عاملیت» را به خمینی نسبت می‌دهد و نقش بازیگران دیگر (بهشتی، رفسنجانی، حزب جمهوری اسلامی) را کم‌اهمیت جلوه می‌دهد. ممکن است خمینی خود اهل فریب نبوده اما توسط حلقهٔ اطرافش در مسیر خاصی سوق داده شده باشد.
    \item مهندسی مجلس خبرگان ممکن است بیشتر کار بهشتی بوده باشد تا خود خمینی.\footnote{رفسنجانی در خاطراتش اشاره می‌کند که بهشتی «معمار» قانون اساسی بود.}
\end{itemize}

% ────────────────────────────────────────
\section{شواهد مکمل: بُعد بین‌المللی}
% ────────────────────────────────────────

مدافعان تز دوم به ابعاد بین‌المللی نیز استناد می‌کنند:

\begin{itemize}[label=\textcolor{cPrimary}{$\blacktriangleright$}, rightmargin=2em]
    \item \textbf{نقش فرانسه:} دولت فرانسه به خمینی در نوفل‌لوشاتو پناه داد و تصویر یک «رهبر معنوی معتدل» از او ساخت. خمینی این تصویر را به سود خود مدیریت کرد.\footnote{\en{Charles Kurzman, The Unthinkable Revolution in Iran, Harvard University Press, 2004, pp.\,67--89.}}
    
    \item \textbf{مصاحبه‌های هدفمند:} خمینی خبرنگاران خاصی را می‌پذیرفت و پاسخ‌ها از قبل با مشاوران (از جمله صادق قطب‌زاده و ابراهیم یزدی) هماهنگ می‌شد. یزدی به‌عنوان مترجم نقش کلیدی داشت و گاهی ترجمه‌هایش «ملایم‌تر» از اصل فارسی بود.\footnote{\en{Moin (1999), pp.\,185--188.} همچنین مقایسه کنید متن فارسی و انگلیسی مصاحبه‌ها در صحیفهٔ امام.}
    
    \item \textbf{پیام به ارتش:} خمینی پیام‌های متعددی به نظامیان فرستاد و قول داد که ارتش حفظ خواهد شد. اما پس از انقلاب، پاکسازی گستردهٔ ارتش (اعدام ده‌ها ژنرال و افسر) نشان داد این وعده‌ها نیز تاکتیکی بودند.\footnote{\en{Bakhash (1984), pp.\,56--72.}}
\end{itemize}

% ────────────────────────────────────────
\section{جمع‌بندی فصل سوم}
% ────────────────────────────────────────

\begin{principleBox}[خلاصهٔ فصل ۳]
\begin{itemize}[label=\textcolor{cPrimary}{$\blacksquare$}, rightmargin=2em]
    \item تز خدعهٔ آگاهانه بر پایهٔ مفاهیم فقهی تقیه و خدعه، نظریهٔ پوپولیسم، و تحلیل ماکیاولیستی استوار است.
    \item قوی‌ترین شواهد: کتاب «حکومت اسلامی» (۱۳۴۸)، عبارت «خدعه» از زبان خمینی، الگوی حذف سیستماتیک، مهندسی قانون اساسی.
    \item مدافعان اصلی: آبراهامیان، گنجی، کاتوزیان، بنی‌صدر.
    \item نقدهای جدی: اتهام توطئه‌اندیشی، نادیده‌گرفتن نقش دیگر بازیگران، ساده‌سازی فرایند تکامل فکری.
    \item \keyword{نقطهٔ قوت اصلی:} توضیح فاصلهٔ عمیق میان وعده و عمل.
    \item \keyword{نقطهٔ ضعف اصلی:} نیاز به فرض یک «طراح کامل» — که با آشفتگی واقعی انقلاب‌ها سازگاری کامل ندارد.
\end{itemize}
\end{principleBox}


% ============================================================
%  فصل ۴: تز سوم — تکامل تدریجی اندیشهٔ ولایت فقیه
% ============================================================
\chapter{تز سوم: تکامل تدریجی اندیشهٔ ولایت فقیه}

\begin{infoBox}[چکیدهٔ فصل]
بر اساس این تز، نه صداقتِ محض و نه فریبِ محض، بلکه \keyword{تکامل تدریجی} اندیشهٔ سیاسی خمینی کلید فهم فاصلهٔ وعده‌ها و عمل است. خمینی در مقاطع مختلف زندگی فکری‌اش مواضع متفاوتی داشت: از «نظارت فقها» در «کشف‌الاسرار» (۱۳۲۲) تا «ولایت فقیه» در درس‌های نجف (۱۳۴۸)، تا «ولایت مطلقهٔ فقیه» در اواخر عمر (۱۳۶۶). وعده‌های پاریس ممکن است بازتاب مرحله‌ای میانی از این تکامل باشند — مرحله‌ای که خمینی هنوز به نقطهٔ نهایی نرسیده بود.
\end{infoBox}

% ────────────────────────────────────────
\section{سیر تحول فکری خمینی: پنج مرحله}
% ────────────────────────────────────────

\begin{figure}[H]
\centering
\begin{tikzpicture}[
    x=3cm,
    stage/.style={
        rectangle, rounded corners=4pt,
        draw=cGold!80, fill=cGold!8,
        text width=4.2cm, align=center,
        font=\small, inner sep=6pt,
        minimum height=3cm
    },
    arrow/.style={
        -{Stealth[length=5pt]},
        line width=1.5pt,
        cGold!60
    },
    label/.style={
        font=\footnotesize\bfseries,
        color=cPrimary
    }
]

\node[stage] (s1) at (0,0) {
    \textbf{مرحله ۱}\\[3pt]
    {\footnotesize ۱۳۲۲ (کشف‌الاسرار)}\\[3pt]
    نظارت فقها بر قوانین\\
    مشروطه + نظارت شرعی\\
    \textcolor{cGreen}{≈ قانون اساسی مشروطه}
};

\node[stage] (s2) at (1.7,0) {
    \textbf{مرحله ۲}\\[3pt]
    {\footnotesize ۱۳۴۱–۱۳۴۳ (قیام ۱۵ خرداد)}\\[3pt]
    مخالفت فعال سیاسی\\
    اما بدون نظریهٔ حکومتی\\
    \textcolor{cWarn}{≈ اپوزیسیون مذهبی}
};

\node[stage] (s3) at (3.4,0) {
    \textbf{مرحله ۳}\\[3pt]
    {\footnotesize ۱۳۴۸ (حکومت اسلامی)}\\[3pt]
    ولایت فقیه = حق حکومت فقها\\
    اما هنوز \keyword{نظری}\\
    \textcolor{cGold}{≈ آرمان‌شهر فقهی}
};

\node[stage] (s4) at (5.1,0) {
    \textbf{مرحله ۴}\\[3pt]
    {\footnotesize ۱۳۵۷ (پاریس)}\\[3pt]
    وعدهٔ عدم حکومت مستقیم\\
    ابهام میان «ناظر» و «حاکم»\\
    \textcolor{cPrimary}{≈ مرحلهٔ گذار}
};

\node[stage] (s5) at (6.8,0) {
    \textbf{مرحله ۵}\\[3pt]
    {\footnotesize ۱۳۶۶ (ولایت مطلقه)}\\[3pt]
    فقیه بالاتر از قانون اساسی\\
    اختیارات نامحدود\\
    \textcolor{cAccent}{≈ ولایت مطلقه}
};

\draw[arrow] (s1) -- (s2);
\draw[arrow] (s2) -- (s3);
\draw[arrow] (s3) -- (s4);
\draw[arrow] (s4) -- (s5);

\end{tikzpicture}
\caption{پنج مرحلهٔ تکامل اندیشهٔ سیاسی خمینی (تز سوم)}
\label{fig:evolution-stages}
\end{figure}

\subsection{مرحلهٔ اول: کشف‌الاسرار (۱۳۲۲/۱۹۴۴)}

نخستین اثر سیاسی خمینی، کتاب «کشف‌الاسرار» بود — پاسخی به کتاب «اسرار هزارساله» نوشتهٔ احمد کسروی (یا منتسب به شاگردان او). در این کتاب، خمینی:

\begin{itemize}[label=\textcolor{cGold}{$\bullet$}, rightmargin=2em]
    \item از نظام \keyword{مشروطه} دفاع می‌کند — مشروط به آنکه قوانین با شرع مغایر نباشد.\footnote{ونسا مارتین تحلیل دقیقی از این متن ارائه داده: \en{Martin (1993), ``Religion and State in Khumayni's Kashf al-Asrar,'' BSOAS, 56(1), 34--45.}}
    \item خواستار \keyword{نظارت} فقها بر قوانین مجلس است — نه حکومت مستقیم آن‌ها.
    \item صراحتاً می‌نویسد: «ما نمی‌گوییم حکومت باید دست فقیه باشد.»\footnote{روح‌الله خمینی، \textit{کشف‌الاسرار}، ص ۱۸۵–۱۸۶ (نسخهٔ چاپ قم).}
    \item در عوض، خواستار آن است که مجلسی از فقها بر قوانین «نظارت شرعی» داشته باشد — مشابه آنچه در قانون اساسی مشروطه (۱۲۸۵) پیش‌بینی شده بود.
\end{itemize}

\begin{noteBox}[نکتهٔ کلیدی]
فاصلهٔ میان «کشف‌الاسرار» (۱۳۲۲) و «حکومت اسلامی» (۱۳۴۸) \keyword{بیست‌وشش سال} است. در این مدت، خمینی از «نظارت فقها در چارچوب مشروطه» به «حکومت مستقیم فقیه» رسید. مدافعان تز سوم استدلال می‌کنند که این تکامل ادامه داشت و در ۱۳۵۷ (پاریس) هنوز به نقطهٔ نهایی (ولایت مطلقه) نرسیده بود.
\end{noteBox}

\subsection{مرحلهٔ دوم: قیام ۱۵ خرداد و تبعید (۱۳۴۱–۱۳۴۳)}

مخالفت خمینی با لایحهٔ انجمن‌های ایالتی و ولایتی (۱۳۴۱) و سپس با «انقلاب سفید» شاه، او را وارد عرصهٔ سیاست فعال کرد. اما در این مرحله:

\begin{itemize}[label=\textcolor{cPrimary}{$\blacktriangleright$}, rightmargin=2em]
    \item خمینی هنوز نظریهٔ حکومتی مدوّنی نداشت.
    \item اعتراضاتش عمدتاً \keyword{سلبی} بود (ضد شاه، ضد وابستگی به آمریکا) نه \keyword{ایجابی} (چه نوع حکومتی باید جایگزین شود).
    \item زبان او بیشتر \keyword{اخلاقی-اعتراضی} بود تا \keyword{نظریه‌پردازانه}.
    \item تبعید به ترکیه (۱۳۴۳) و سپس نجف (۱۳۴۳–۱۳۵۷) فرصت تأمل نظری به او داد.\footnote{\en{Moin (1999), pp.\,86--120.}}
\end{itemize}

\subsection{مرحلهٔ سوم: درس‌های نجف و «حکومت اسلامی» (۱۳۴۸)}

در زمستان ۱۳۴۸/۱۹۷۰، خمینی سلسله درس‌هایی در حوزهٔ علمیهٔ نجف ارائه داد که بعدها با عنوان «ولایت فقیه» یا «حکومت اسلامی» منتشر شد. این متن جهشی کیفی (\en{qualitative leap}) در اندیشهٔ سیاسی اوست:

\begin{warningBox}[تحول کلیدی: از «نظارت» به «ولایت»]
\begin{longtable}{%
    >{\raggedleft\arraybackslash}p{0.20\textwidth}
    >{\raggedleft\arraybackslash}p{0.35\textwidth}
    >{\raggedleft\arraybackslash}p{0.35\textwidth}}
\toprule
\rowcolor{cAccent!10}
\textbf{مفهوم} & \textbf{کشف‌الاسرار (۱۳۲۲)} & \textbf{حکومت اسلامی (۱۳۴۸)} \\
\midrule

نقش فقیه & ناظر بر قوانین & \keyword{حاکم} و دارای ولایت \\
\midrule
دامنهٔ اختیارات & محدود به تطبیق شرعی & شامل تمام امور حکومتی \\
\midrule
مبنای مشروعیت & قانون اساسی مشروطه & \keyword{حکم الهی} مستقل از قانون اساسی \\
\midrule
نسبت با دموکراسی & سازگار با مشروطه & مقدم بر رأی مردم \\
\midrule
سلطنت & قابل اصلاح & باید \keyword{سرنگون} شود \\

\bottomrule
\end{longtable}
\end{warningBox}

مدافعان تز سوم (به‌ویژه ارجمند و مارتین) تأکید می‌کنند که حتی در «حکومت اسلامی»، نظریه هنوز \keyword{انتزاعی} و \keyword{آرمانی} بود. خمینی هیچ نقشهٔ عملی برای اجرای آن ارائه نداد — مثلاً دربارهٔ ساختار حکومت، نحوهٔ انتخاب فقیه، یا رابطهٔ ولی فقیه با مجلس سخنی نگفت.\footnote{\en{Arjomand (1988), pp.\,147--155.}}

\subsection{مرحلهٔ چهارم: پاریس (۱۳۵۷) — مرحلهٔ ابهام}

مدافعان تز سوم استدلال می‌کنند که دورهٔ پاریس \keyword{مرحلهٔ گذار} در فکر خمینی بود:

\begin{principleBox}[تحلیل مرحلهٔ پاریس — تز سوم]
\begin{itemize}[label=\textcolor{cGold}{$\bullet$}, rightmargin=2em]
    \item خمینی در ۱۳۴۸ نظریه‌ای \textit{داشت} اما هنوز مطمئن نبود آیا عملاً قابل اجراست یا نه.
    \item در پاریس، تحت تأثیر مشاوران لیبرال‌تر (یزدی، قطب‌زاده، بازرگان)، ممکن است واقعاً \keyword{تردید} کرده باشد.
    \item وعدهٔ «عدم حکومت» نه لزوماً «دروغ» بود و نه لزوماً «حقیقت کامل» — بلکه بازتاب \keyword{عدم قطعیت} فکری بود.
    \item خمینی هنوز نمی‌دانست نظام آینده دقیقاً چه شکلی خواهد داشت.
    \item فقط پس از بازگشت به ایران و مواجهه با واقعیت‌های قدرت، اندیشه‌اش به نقطهٔ نهایی رسید.
\end{itemize}
\end{principleBox}

\subsection{مرحلهٔ پنجم: ولایت مطلقهٔ فقیه (۱۳۶۶)}

آخرین و رادیکال‌ترین تحول فکری خمینی در دی‌ماه ۱۳۶۶ رخ داد — زمانی که در نامه‌ای به خامنه‌ای (رئیس‌جمهور وقت) اعلام کرد:

\begin{warningBox}[نقل‌قول — نامهٔ دی ۱۳۶۶]
\textit{«حکومت، شعبه‌ای از ولایت مطلقهٔ رسول‌الله (ص) است... حکومت می‌تواند از احکام اولیهٔ اسلام هم جلوگیری کند... حکومت می‌تواند هر امری را — حتی نماز و روزه و حج را — موقتاً تعطیل کند.»}

\hfill — آیت‌الله خمینی، نامه به خامنه‌ای، ۱۶ دی ۱۳۶۶\footnote{صحیفهٔ امام، ج۲۰، ص ۴۵۱–۴۵۲.}
\end{warningBox}

این اظهارنظر \keyword{بی‌سابقه} بود — حتی در «حکومت اسلامی» (۱۳۴۸) چنین ادعایی نشده بود. مدافعان تز سوم از این نامه به‌عنوان شاهدی بر \keyword{تکامل مستمر} بهره می‌گیرند: اگر خمینی حتی بین ۱۳۵۸ (قانون اساسی) و ۱۳۶۶ (ولایت مطلقه) تحول فکری داشت، چرا نباید باور کنیم بین ۱۳۴۸ و ۱۳۵۷ هم تحول داشته باشد؟

% ────────────────────────────────────────
\section{مدافعان اصلی تز سوم}
% ────────────────────────────────────────

\subsection{سعید امیر ارجمند}

ارجمند در \en{The Turban for the Crown} (۱۹۸۸) یکی از دقیق‌ترین تحلیل‌ها را از تکامل نهادی ولایت فقیه ارائه می‌دهد. نکات کلیدی او:

\begin{itemize}[label=\textcolor{cGold}{$\blacktriangleright$}, rightmargin=2em]
    \item نظریهٔ ولایت فقیه در نجف یک «ایدهٔ خام» بود. تبدیل آن به یک \keyword{نهاد حقوقی} کار مجلس خبرگان (به‌ویژه بهشتی) بود، نه خود خمینی.\footnote{\en{Arjomand (1988), pp.\,150--168.}}
    \item خمینی در پاریس ممکن بود واقعاً تصور کند مدل نهایی چیزی میان «نظارت» و «ولایت» خواهد بود — و مجلس خبرگان این‌را به سمت «ولایت» تمام‌عیار کشاند.
    \item تصمیم خمینی به تأسیس مجلس خبرگان (به‌جای رفراندوم مستقیم بر پیش‌نویس) لحظهٔ تعیین‌کننده بود — اما ممکن بود این تصمیم خودش محصول تأثیرپذیری از حلقهٔ درونی باشد.
\end{itemize}

\subsection{ونسا مارتین}

مارتین در \en{Creating an Islamic State} (۲۰۰۰) تحلیلی متنی (\en{textual analysis}) ارائه می‌دهد:

\begin{noteBox}[یافته‌های مارتین]
\begin{enumerate}[label=\textcolor{cGold}{\bfseries\arabic*.}, rightmargin=2em]
    \item مقایسهٔ دقیق «کشف‌الاسرار» (۱۳۲۲) و «حکومت اسلامی» (۱۳۴۸) نشان‌دهندهٔ \keyword{تحول واقعی} است — نه صرفاً تفاوت مخاطب.
    \item اصطلاحات فقهی خمینی در دو متن \keyword{ساختاراً متفاوت} است: در ۱۳۲۲ از «نظارت» (\rl{إشراف}) و در ۱۳۴۸ از «ولایت» (\rl{ولایة}) استفاده شده.
    \item این تغییر اصطلاحی بازتاب تغییر \keyword{مبنای فقهی} است: از فقه مشروطه‌محور نائینی به فقه حکومتی مستقل.\footnote{\en{Martin (2000), pp.\,94--120.}}
    \item مصاحبه‌های پاریس از نظر اصطلاحی بیشتر به «کشف‌الاسرار» نزدیکند تا «حکومت اسلامی» — که ممکن است نشانهٔ \keyword{عقب‌نشینی موقت} فکری باشد.
\end{enumerate}
\end{noteBox}

\subsection{محسن کدیور}

کدیور از زاویهٔ \keyword{فقه تطبیقی} تز سوم را تقویت می‌کند:

\begin{itemize}[label=\textcolor{cPrimary}{$\blacktriangleright$}, rightmargin=2em]
    \item ولایت فقیه در تاریخ فقه شیعه \keyword{اجماعی نبوده} و مراجع بزرگ (بروجردی، خویی، شریعتمداری) آن‌را نپذیرفتند.\footnote{محسن کدیور، \textit{نظریه‌های دولت در فقه شیعه}، فصل ۶.}
    \item بنابراین، حتی خود خمینی ممکن بود در \keyword{مشروعیت فقهی} نظریه‌اش تردید داشته باشد — حداقل تا زمانی که قدرت عملی نظریه را «اثبات» کرد.
    \item «ولایت مطلقه» (۱۳۶۶) خود اعترافی ضمنی است: اگر خمینی از ابتدا ولایت مطلقه می‌خواست، چرا قانون اساسی ۱۳۵۸ آن‌را \keyword{نداشت}؟ چرا بازنگری ۱۳۶۸ لازم شد؟\footnote{قانون اساسی ۱۳۵۸ ولایت فقیه را داشت اما «مطلقه» نبود. اصلاحیهٔ ۱۳۶۸ واژهٔ «مطلقه» را اضافه کرد.}
\end{itemize}

% ────────────────────────────────────────
\section{نقد تز سوم}
% ────────────────────────────────────────

\subsection{نقد از منظر تز دوم (خدعه)}

\begin{itemize}[label=\textcolor{cAccent}{$\times$}, rightmargin=2em]
    \item \textbf{«تکامل» یا «پنهان‌کاری»؟} آبراهامیان استدلال می‌کند که آنچه تز سوم «تکامل فکری» می‌نامد، ممکن است صرفاً \keyword{پنهان‌سازی تاکتیکی} باشد. خمینی در ۱۳۴۸ نظریه‌اش را نزد شاگردان حوزوی گفت (که مخاطبان محدودی بودند) و در ۱۳۵۷ آن‌را نزد رسانه‌های جهانی پنهان کرد.\footnote{\en{Abrahamian (1993), pp.\,23--25.}}
    \item \textbf{عدم شاهد بر «تردید»:} هیچ سند مستقیمی وجود ندارد که نشان دهد خمینی در پاریس واقعاً دربارهٔ ولایت فقیه \keyword{تردید} کرد. ادعای «تردید» یک \keyword{حدس} است، نه یک یافته.
    \item \textbf{مسئلهٔ «اختیار اوکام»:} (\en{Occam's Razor}) اگر خمینی هم نظریه داشت (۱۳۴۸) و هم آن‌را اجرا کرد (۱۳۵۸ به بعد)، ساده‌ترین توضیح آن است که از ابتدا قصد اجرا داشت — نه اینکه تکامل پیچیده‌ای طی کرد.
\end{itemize}

\subsection{نقد از منظر تز اول (صداقت)}

\begin{itemize}[label=\textcolor{cGreen}{$\times$}, rightmargin=2em]
    \item تز سوم هنوز مسئلهٔ «اجبار بحران‌ها» را جدی نمی‌گیرد. حتی اگر تکامل فکری وجود داشت، \keyword{سرعت} تمرکز قدرت بیشتر با فشار بحران‌ها توضیح‌پذیر است تا با تحول فکری.
\end{itemize}

\subsection{نقد درونی}

\begin{itemize}[label=\textcolor{cGold}{$\times$}, rightmargin=2em]
    \item تز سوم بیشتر یک \keyword{چارچوب تبیینی} است تا یک «تز» مستقل. می‌تواند با هر یک از تزهای دیگر ترکیب شود: «صادق بود + تکامل یافت» یا «فریب داد + تکامل هم یافت». به‌همین‌دلیل برخی پژوهشگران آن‌را نه یک تز رقیب بلکه یک \keyword{مکمل} تلقی می‌کنند.
\end{itemize}

% ────────────────────────────────────────
\section{جدول تطبیقی: مراحل تکامل و مواضع}
% ────────────────────────────────────────

\begin{longtable}{%
    >{\centering\arraybackslash}p{0.08\textwidth}
    >{\raggedleft\arraybackslash}p{0.12\textwidth}
    >{\raggedleft\arraybackslash}p{0.22\textwidth}
    >{\raggedleft\arraybackslash}p{0.22\textwidth}
    >{\raggedleft\arraybackslash}p{0.22\textwidth}}
\toprule
\rowcolor{cGold!15}
\textbf{مرحله} & \textbf{سال} & \textbf{متن کلیدی} & \textbf{موضع دربارهٔ حکومت} & \textbf{دامنهٔ اختیارات فقیه} \\
\midrule
\endhead

\cellcolor{cGreen!8} ۱ & \cellcolor{cGreen!8} ۱۳۲۲ & \cellcolor{cGreen!8} کشف‌الاسرار & \cellcolor{cGreen!8} نظارت بر قوانین، حفظ مشروطه & \cellcolor{cGreen!8} محدود به تطبیق شرعی \\
\midrule

\cellcolor{cGold!8} ۲ & \cellcolor{cGold!8} ۱۳۴۱–۴۳ & \cellcolor{cGold!8} سخنرانی‌ها و اعلامیه‌ها & \cellcolor{cGold!8} ضدسلطنت، بدون نظریهٔ جایگزین & \cellcolor{cGold!8} نامشخص \\
\midrule

\cellcolor{cWarn!8} ۳ & \cellcolor{cWarn!8} ۱۳۴۸ & \cellcolor{cWarn!8} حکومت اسلامی & \cellcolor{cWarn!8} حکومت مستقیم فقیه & \cellcolor{cWarn!8} ولایت فقیه (نظری و انتزاعی) \\
\midrule

\cellcolor{cPrimary!8} ۴ & \cellcolor{cPrimary!8} ۱۳۵۷ & \cellcolor{cPrimary!8} مصاحبه‌های پاریس & \cellcolor{cPrimary!8} «مقامی نخواهم داشت» & \cellcolor{cPrimary!8} ابهام‌آمیز \\
\midrule

\cellcolor{cAccent!8} ۵ & \cellcolor{cAccent!8} ۱۳۶۶ & \cellcolor{cAccent!8} نامه به خامنه‌ای & \cellcolor{cAccent!8} ولایت مطلقه بالاتر از احکام اولیه & \cellcolor{cAccent!8} نامحدود \\

\bottomrule
\caption{مراحل تکامل اندیشهٔ سیاسی خمینی — جدول تطبیقی}
\label{tab:evolution-stages}
\end{longtable}

% ────────────────────────────────────────
\section{جمع‌بندی فصل چهارم}
% ────────────────────────────────────────

\begin{principleBox}[خلاصهٔ فصل ۴]
\begin{itemize}[label=\textcolor{cPrimary}{$\blacksquare$}, rightmargin=2em]
    \item تز تکامل تدریجی بر پایهٔ تحلیل متنی آثار خمینی در مقاطع مختلف استوار است.
    \item پنج مرحلهٔ تکامل قابل شناسایی: نظارت (۱۳۲۲) → مخالفت سیاسی (۱۳۴۱) → ولایت نظری (۱۳۴۸) → ابهام (۱۳۵۷) → ولایت مطلقه (۱۳۶۶).
    \item مدافعان اصلی: ارجمند، مارتین، کدیور.
    \item \keyword{نقطهٔ قوت:} توضیح تغییرات مستمر در مواضع خمینی بدون نیاز به فرض «فریب» یا «اجبار».
    \item \keyword{نقطهٔ ضعف:} ابهام ذاتی (هر تفسیری ممکن است)، فقدان شاهد مستقیم بر «تردید» خمینی، قابلیت ترکیب با تزهای دیگر (استقلال ضعیف).
\end{itemize}
\end{principleBox}


% ============================================================
%  فصل ۵: تز چهارم — فشار نخبگان و دینامیک‌های گروهی
% ============================================================
\chapter{تز چهارم: فشار نخبگان و دینامیک‌های گروهی}

\begin{infoBox}[چکیدهٔ فصل]
بر اساس این تز، تمرکز صرف بر شخص خمینی — چه به‌عنوان فریبکار و چه به‌عنوان قربانی شرایط — \keyword{ناکافی} است. کلید فهم تمرکز قدرت، \keyword{حلقهٔ نخبگان انقلابی} است: شبکه‌ای از روحانیون سیاسی، تکنوکرات‌های مذهبی، و فعالان انقلابی که منافع مشترک در تمرکز قدرت داشتند و خمینی را — آگاهانه یا ناآگاهانه — در مسیر رهبری تمام‌عیار سوق دادند. خمینی نه فریبکار بود و نه قربانی؛ بلکه \keyword{بازیگر اصلی در یک بازی گروهی} بود.
\end{infoBox}

% ────────────────────────────────────────
\section{بنیان‌های نظری تز چهارم}
% ────────────────────────────────────────

\subsection{نظریهٔ رقابت نخبگان}

تز چهارم از چارچوب‌های نظری \keyword{جامعه‌شناسی سیاسی} بهره می‌گیرد، به‌ویژه:

\begin{itemize}[label=\textcolor{cConsolid}{$\bullet$}, rightmargin=2em]
    \item \textbf{گائتانو موسکا} (\en{Gaetano Mosca}) و \textbf{ویلفردو پارتو} (\en{Vilfredo Pareto}): هر نظام سیاسی توسط «طبقهٔ حاکم» (نخبگان) اداره می‌شود. انقلاب‌ها نه حذف نخبگان بلکه \keyword{جایگزینی} یک طبقهٔ نخبه با طبقه‌ای دیگر هستند.\footnote{\en{Vilfredo Pareto, The Rise and Fall of the Elites (1901); Gaetano Mosca, The Ruling Class (1939).}}
    
    \item \textbf{رابرت میشلز} (\en{Robert Michels}): «قانون آهنین الیگارشی» — هر سازمان، هرچقدر هم دموکراتیک شروع کند، به‌سمت تمرکز قدرت در دست عدهٔ کمی حرکت می‌کند.\footnote{\en{Robert Michels, Political Parties: A Sociological Study of the Oligarchical Tendencies of Modern Democracy (1911).}}
    
    \item \textbf{تدا اسکاچپول} (\en{Theda Skocpol}): انقلاب‌ها محصول فروپاشی ساختار دولت هستند، نه لزوماً طراحی یک فرد. \keyword{خلأ قدرت} فضا را برای رقابت نخبگان باز می‌کند.\footnote{\en{Theda Skocpol, States and Social Revolutions (1979).}}
\end{itemize}

\subsection{مفهوم «حلقهٔ درونی» (\en{Inner Circle})}

\begin{principleBox}[حلقهٔ درونی خمینی — بازیگران کلیدی]
\begin{longtable}{%
    >{\raggedleft\arraybackslash}p{0.18\textwidth}
    >{\raggedleft\arraybackslash}p{0.22\textwidth}
    >{\raggedleft\arraybackslash}p{0.25\textwidth}
    >{\raggedleft\arraybackslash}p{0.25\textwidth}}
\toprule
\rowcolor{cConsolid!10}
\textbf{شخص} & \textbf{نقش} & \textbf{منفعت در تمرکز قدرت} & \textbf{سرنوشت} \\
\midrule
\endhead

محمد بهشتی &
معمار قانون اساسی، رئیس دیوان عالی &
تثبیت ولایت فقیه = تضمین قدرت روحانیت &
ترور ۷ تیر ۱۳۶۰ \\
\midrule

اکبر هاشمی رفسنجانی &
رئیس مجلس، فرمانده جنگ &
تمرکز قدرت = حذف رقبای غیرروحانی &
رئیس‌جمهور ۱۳۶۸–۱۳۷۶ \\
\midrule

علی خامنه‌ای &
رئیس‌جمهور ۱۳۶۰–۱۳۶۸ &
ولایت فقیه = چارچوب نهادی برای جانشینی &
ولی فقیه از ۱۳۶۸ \\
\midrule

محمدجواد باهنر &
نخست‌وزیر (کوتاه‌مدت) &
حکومت اسلامی = تحقق آرمان‌ها &
ترور ۸ شهریور ۱۳۶۰ \\
\midrule

احمد خمینی &
واسطهٔ پدر با جهان بیرون &
تداوم «خانوادهٔ خمینی» در مرکز قدرت &
مرگ مشکوک ۱۳۷۳ \\
\midrule

صادق خلخالی &
حاکم شرع دادگاه‌های انقلاب &
قدرت قضایی بدون نظارت &
بازنشستگی اجباری \\
\midrule

محمدعلی رجایی &
نخست‌وزیر، رئیس‌جمهور &
نمایندهٔ جناح تندرو غیرروحانی &
ترور ۸ شهریور ۱۳۶۰ \\

\bottomrule
\end{longtable}
\end{principleBox}

% ────────────────────────────────────────
\section{مکانیسم‌های فشار نخبگان}
% ────────────────────────────────────────

مدافعان تز چهارم چند مکانیسم مشخص شناسایی می‌کنند که از طریق آن‌ها نخبگان انقلابی خمینی را به‌سمت تمرکز قدرت سوق دادند:

\subsection{مکانیسم ۱: مهندسی مجلس خبرگان قانون اساسی}

\begin{warningBox}[نقش بهشتی در مهندسی مجلس خبرگان]
آیت‌الله محمد بهشتی، دبیرکل حزب جمهوری اسلامی، نقش محوری در سه فرایند داشت:

\begin{enumerate}[label=\textcolor{cAccent}{\bfseries\arabic*.}, rightmargin=2em]
    \item \textbf{اقناع خمینی} برای تشکیل مجلس خبرگان به‌جای رفراندوم مستقیم بر پیش‌نویس.
    \item \textbf{سازمان‌دهی} فهرست نامزدهای حزب جمهوری اسلامی برای انتخابات مجلس خبرگان — که به پیروزی قاطع روحانیون تندرو منجر شد.\footnote{از ۷۳ عضو مجلس خبرگان، حدود ۵۵ نفر روحانی بودند و اکثر آن‌ها وابسته به حزب جمهوری اسلامی. نگاه کنید به: \en{Schirazi (1997), pp.\,30--38.}}
    \item \textbf{مدیریت مذاکرات} مجلس خبرگان به‌عنوان نایب‌رئیس — و هدایت بحث‌ها به‌سمت ولایت فقیه.
\end{enumerate}
\end{warningBox}

مدافعان تز چهارم تأکید می‌کنند که بهشتی — نه خمینی — «معمار اصلی» ولایت فقیه در قانون اساسی بود. خمینی نقش «تأییدکننده» داشت، نه «طراح».\footnote{رفسنجانی در خاطراتش نقل می‌کند: «آقای بهشتی بار اصلی کار [قانون اساسی] را بر دوش داشت. امام تأیید می‌فرمودند.» هاشمی رفسنجانی، \textit{عبور از بحران}، ص ۱۸۵.}

\subsection{مکانیسم ۲: ایجاد بحران و بهره‌برداری از آن}

برخی مدافعان تز چهارم (با نزدیکی به تز دوم) استدلال می‌کنند که حلقهٔ درونی نه‌تنها از بحران‌ها بهره‌برداری کرد بلکه گاهی آن‌ها را \keyword{دامن زد}:

\begin{itemize}